\subsection{Comparison with the float-INN baseline}
Table~\ref{tab:compare} summarises all methods at 1024 bits. The float-INN baseline hides the message well (its decoder
makes no errors and its stego has \pinnBASE{}~dB before compensation), but when the receiver inverts it with the latent
set to zero and the stego rounded to 8 bits, \resBASE\% of the cover pixels come back wrong. Coding that residual needs
\auxBASE{} bits on average, about \auxRatio{} times more than IC-INN-Q and far more than the
\capPEE{} bits a single PEE layer offers on these images, so not one of the \nImgs{} images could be embedded
reversibly. This is exactly the failure that motivated the integer, message-conditioned design: in IC-INN the cover
residual is zero by construction and only overflow and decoder errors remain.

\begin{table*}[t]
\centering
\caption{Comparison at 1024 bits (0.0156 bpp) over all test images. Success = embedding possible; Exact = cover and message recovered bit for bit; lossy-channel BER for Gaussian noise ($\sigma{=}2$) and 1\% salt-and-pepper noise. Times are MATLAB/Octave on a 2-core CPU. $^\dagger$PSNR of the INN stego before PEE; its side information never fits.}
\label{tab:compare}
\setlength{\tabcolsep}{4pt}
\begin{tabular}{lccccccccc}
\toprule
Method & Params & PSNR (dB) & SSIM & Aux (bits) & Success & Exact & BER GN2 (\%) & BER S\&P1 (\%) & Embed / extract (s)\\
\midrule
Rhombus PEE & 0 & 65.32 & 0.9998 & -- & 132/132 & 132/132 & 49.87 & 48.99 & 0.02 / 0.02\\
Float INN + RC (baseline) & 128.7k & 46.07$^\dagger$ & -- & 111.7k & 0/132 & 0/132 & -- & -- & -- \\
IC-INN-Q (ours) & 62.9k & 49.05 & 0.9949 & 561 & 131/132 & 131/132 & 3.22 & 6.10 & 0.51 / 0.51\\
IC-INN-R (ours) & 62.9k & 47.57 & 0.9940 & 610 & 128/132 & 128/132 & 1.76 & 4.03 & 0.51 / 0.51\\
\bottomrule
\end{tabular}
\end{table*}


\subsection{Robustness on a lossy channel}\label{sec:robust}
Table~\ref{tab:robust} reports the BER when the final stego is distorted before extraction (all Kodak and DICOM images
and the first 18 BSD100 images). Classical PEE loses everything: its extraction depends on exact pixel values, so its BER
is about 50\% under every attack. Our decoder reads the message directly from the received image. Without any attack it
already gets \robnonezeroR\% of the bits wrong (IC-INN-R), which is why the side information is needed for exact recovery;
under Gaussian noise with $\sigma=2$ the BER is \robgausstwoR\% and under 1\% salt-and-pepper noise \robspsponeR\%. Mild
JPEG compression (quality 95) gives \robjpegnfR\%, but quality 90 or lower destroys the message (\robjpegninetyR\%), which is
not surprising since JPEG was not part of the noise layer. The robustness-oriented model is consistently better than the
quality model, at the cost of 1.5~dB.

\begin{table}[t]
\centering
\caption{Message BER (\%) on a lossy channel, payload 1024 bits, 50 images. The cover is not recoverable here; the question is only how many message bits survive.}
\label{tab:robust}
\setlength{\tabcolsep}{4pt}
\begin{tabular}{lcccc}
\toprule
Attack & PEE & Float INN & IC-INN-Q & IC-INN-R\\
\midrule
none & 0.00 & -- & 1.78 & 1.16\\
GN $\sigma{=}1$ & 50.17 & -- & 2.07 & 1.26\\
GN $\sigma{=}2$ & 49.87 & -- & 3.22 & 1.76\\
GN $\sigma{=}3$ & 50.05 & -- & 6.30 & 3.62\\
S\&P 0.5\% & 48.29 & -- & 3.35 & 2.17\\
S\&P 1\% & 48.99 & -- & 6.10 & 4.03\\
JPEG 95 & 50.15 & -- & 12.14 & 5.07\\
JPEG 90 & 50.47 & -- & 38.74 & 36.07\\
JPEG 80 & 50.04 & -- & 47.02 & 46.10\\
\bottomrule
\end{tabular}
\end{table}


\subsection{Complexity}
The coupling networks have 42{,}888 parameters and the decoder 19{,}969. Embedding and extraction take about
\tembQfour{}~s each for a $256\times256$ image in our MATLAB code (Octave, single CPU core), of which most is the inverse
pass and the arithmetic coder. Classical PEE takes about 0.02~s. The GUI (\texttt{INNRC\_App.m}) runs the same functions
interactively.
